\section{The Aumkara of The West}

As is well known, the Mandukya Upanishad explicates the individual letters of the sacred monosyllable ``AUM'' as symbols of Atma within various states, and finally with the twelfth verse describes the AUM taken as a whole as the ``state'' of Turiya: the identity of Aum and the Self, the ineffable reality of Atma knowing Atma, for, by, and through Atma.

Recently, in his Living the Dream\footnote{\url{https://www.gornahoor.net/?p=5811}} essay, Cologero summarized the ``description'' of Turiya as found within the Vedanta Tradition, in ``via negativia'' language; with the important conclusion being drawn, that the state need be cultivated by mastering thought, and stilling the mind. As opposed to repeating those ``descriptions'' herein, readers may wish to review the respective entry before proceeding.

Since metaphysics as ``Universal science'' is not the ``property'' of any religious form, as Frithjof Schuon put it, one might then anticipate within the Western Tradition to locate something equivalent or parallel to the elucidation of Aum. And indeed, with the Neo-Platonic ``To Hen'', The One, is such, a rather obvious ineffable equal of Atma or Para Brahman.

In his \textit{Commentary on Plato's Parmenides}, Proclus explains that while there can really be no ``name'' to assign to the Principle he is speaking about, the best that can be done is to settle upon ``hen'', although calling it ``The One'' must never be understood as anything other than a place holder for what can never truly be named. Still, the word ``hen'' itself,

Proclus points out, contains some clues to assist our contemplations---clues moreover worth comparing with the Mandukya:

\begin{enumerate}
\item Proclus states the One is ``beyond breath'', ``But the first of the things that emanate from it is represented by the rough breathing with which we utter `hen'. Itself, it is unnamable, just as the breathing by itself is silent''. The ``A'' of Aum is ``all pervasive'', unmanifest, yet the principle of all manifestation--while well near the Absolute, it still is not quite the metaphysical ``beyond Being''; similarly, the hen cycle proceeds with ``rough breathing''--primitive, the rudiments of all breath that could be breathed, and yet not indistinct enough to be not breath, or ``beyond breath'' (The One).

\item Proclus continues, ``The second is represented by the utterable vowel which now becomes utterable with the breathing, and it in itself becomes both utterable and unutterable, unspeakable and speakable; for the procession of the second order of existence has to be mediated''. Similarly, the ``U'' of Aum has the function of mediator in the Mandukya Upansishad, not just as occupying the central place within the word, but as the dream state (hence between waking and dreamless sleep), and more importantly by analogy, as the ``psychic'' world, the ``Mundus Imaginalis'' betwixt the corporeal and spiritual worlds. The ``lower'' bounds communicating with the corporeal; the ``upper'' bounds communicating with the spiritual, which is why Proclus calls it ``utterable AND unutterable'', ``unspeakable AND speakable'', as it participates in both directions.

\item ``Third comes hen which contains the unsounding breathing and the soundable force of `e' and the letter that goes with this, the consonant `n', which represents in a converse way the same thing as breathing''. As it ``represents in a converse way'' breathing, it is a return to breathing, which began with the ``rough breathing'' of hen (1); similarly, the ``M'' of Aum is in Mudukaya a type of return to the unmanifested in the ``A'', or the ``beyond breath'' at the initiatilization of hen.

\item Proclus concludes this section noting ``The whole is a triad formed in this way: from this one derives a dyad, but behind the dyad there is a monad. But the first principle is beyond everything and not merely beyond this triad which is the first thing that comes after it... .that prior breathing which is the silent symbol of Being''. This whole then, more than just sharing an accord with previous ``descriptions'' of Turiya as in the mentioned Cologero essay, is to be found only in a certain profound, ``prior'' (primordial) ``silence''.
\end{enumerate}

\flrightit{Posted on 2013-02-20 by Frater M}

